%% Prague Cafe Project -- presentation (topic: Ratings)
%% Compile from the folder presentation/. Needs ../report/tables/numbers.tex and ../report/figures/
\documentclass[aspectratio=169,11pt]{beamer}
\usetheme{Madrid}
\usecolortheme{default}
\setbeamertemplate{navigation symbols}{}
\usepackage[utf8]{inputenc}
\usepackage[T1]{fontenc}
\usepackage{booktabs}
\usepackage{graphicx}
\graphicspath{{../report/figures/}}

% same statistics macros as in the report (cafestyle.sty)
\newcommand{\pr}{\textit{p}}
\newcommand{\Ncol}{\textit{N}}
\newcommand{\CI}[2]{95\% CI [#1,\,#2]}
\newcommand{\tstat}[2]{\textit{t}(#1)\,=\,#2}
\newcommand{\Fstat}[3]{\textit{F}(#1,\,#2)\,=\,#3}
\newcommand{\chisqstat}[3]{$\chi^2$(#1, \Ncol\,=\,#2)\,=\,#3}
\newcommand{\pval}[1]{\pr\,=\,#1}
\newcommand{\plt}[1]{\pr\,<\,#1}
\newcommand{\AoneN}{186}
\newcommand{\AoneEst}{4.54}
\newcommand{\AoneLo}{4.51}
\newcommand{\AoneHi}{4.58}
\newcommand{\AoneDiff}{0.04}
\newcommand{\AoneEff}{0.17}
\newcommand{\AoneStat}{\tstat{185}{2.36}}
\newcommand{\AoneP}{\pval{.019}}
\newcommand{\AtwoN}{185}
\newcommand{\AtwoEst}{-0.04}
\newcommand{\AtwoLo}{-0.12}
\newcommand{\AtwoHi}{0.03}
\newcommand{\AtwoEff}{-0.17}
\newcommand{\AtwoEffLo}{-0.46}
\newcommand{\AtwoEffHi}{0.12}
\newcommand{\AtwoStat}{\tstat{176.3}{-1.17}}
\newcommand{\AtwoP}{\pval{.123}}
\newcommand{\AthreeN}{186}
\newcommand{\AthreeEst}{0.010}
\newcommand{\AthreeLo}{-0.003}
\newcommand{\AthreeHi}{0.024}
\newcommand{\AthreeEff}{0.11}
\newcommand{\AthreeEffLo}{-0.03}
\newcommand{\AthreeEffHi}{0.26}
\newcommand{\AthreeStat}{\tstat{185}{1.54}}
\newcommand{\AthreeP}{\pval{.125}}
\newcommand{\AfourN}{186}
\newcommand{\AfourStat}{\Fstat{2}{50.0}{12.05}}
\newcommand{\AfourP}{\plt{.001}}
\newcommand{\AfourEff}{0.15}
\newcommand{\AfourEffLo}{0.06}
\newcommand{\AfourEffHi}{0.31}
\newcommand{\AfiveN}{185}
\newcommand{\AfiveEst}{15.5}
\newcommand{\AfiveLo}{1.8}
\newcommand{\AfiveHi}{29.3}
\newcommand{\AfiveEff}{0.16}
\newcommand{\AfiveEffLo}{0.03}
\newcommand{\AfiveEffHi}{0.30}
\newcommand{\AfiveStat}{\chisqstat{1}{185}{4.80}}
\newcommand{\AfiveP}{\pval{.014}}
\newcommand{\AfiveMinExp}{33.8}
\newcommand{\AsixN}{186}
\newcommand{\AsixEst}{-0.12}
\newcommand{\AsixLo}{-0.27}
\newcommand{\AsixHi}{0.04}
\newcommand{\AsixEff}{0.015}
\newcommand{\AsixP}{\pval{.096}}
\newcommand{\AsixPearson}{-0.09}
\newcommand{\AsixPearsonLo}{-0.23}
\newcommand{\AsixPearsonHi}{0.05}
\newcommand{\AsixPearsonP}{\pval{.206}}
\newcommand{\AsixKendall}{-0.10}
\newcommand{\AsixKendallP}{\pval{.063}}
\newcommand{\AsevenN}{186}
\newcommand{\AsevenEst}{-0.011}
\newcommand{\AsevenLo}{-0.029}
\newcommand{\AsevenHi}{0.006}
\newcommand{\AsevenEff}{0.009}
\newcommand{\AsevenP}{\pval{.206}}
\newcommand{\AsevenBP}{.318}
\newcommand{\AsevenXEst}{-0.029}
\newcommand{\AsevenXLo}{-0.059}
\newcommand{\AsevenXHi}{0.001}
\newcommand{\AsevenXP}{\pval{.055}}
\newcommand{\AeightN}{185}
\newcommand{\AeightStat}{\Fstat{2}{181}{5.04}}
\newcommand{\AeightP}{\pval{.007}}
\newcommand{\AeightIntEst}{-0.287}
\newcommand{\AeightIntLo}{-0.467}
\newcommand{\AeightIntHi}{-0.108}
\newcommand{\AeightLocEst}{-0.007}
\newcommand{\AeightLocLo}{-0.090}
\newcommand{\AeightLocHi}{0.076}
\newcommand{\AeightSlopeSimple}{-0.011}
\newcommand{\AeightSlopeMult}{-0.012}
\newcommand{\AeightNestedF}{\Fstat{2}{181}{5.04}, \pval{.007}}
\newcommand{\AeightVIFmax}{1.02}
\newcommand{\AeightEff}{0.061}
\newcommand{\AnineN}{186}
\newcommand{\AnineEst}{0.497}
\newcommand{\AnineLo}{0.401}
\newcommand{\AnineHi}{0.592}
\newcommand{\AnineEff}{0.983}
\newcommand{\AnineP}{\plt{.001}}
\newcommand{\AnineBP}{\plt{.001}}
\newcommand{\AnineNestedF}{\Fstat{1}{182}{675.82}, \plt{.001}}
\newcommand{\AnineSlopeU}{0.396}
\newcommand{\AnineSlopeULo}{0.304}
\newcommand{\AnineSlopeUHi}{0.487}
\newcommand{\AnineSlopeO}{0.892}
\newcommand{\AnineSlopeOLo}{0.863}
\newcommand{\AnineSlopeOHi}{0.921}
\newcommand{\AtenN}{186}
\newcommand{\AtenEvents}{69}
\newcommand{\AtenEst}{-12.8}
\newcommand{\AtenLo}{-23.2}
\newcommand{\AtenHi}{-5.6}
\newcommand{\AtenOR}{0.56}
\newcommand{\AtenORLo}{0.36}
\newcommand{\AtenORHi}{0.87}
\newcommand{\AtenStat}{\textit{z}\,=\,-2.57}
\newcommand{\AtenP}{\pval{.010}}
\newcommand{\AtenEff}{0.044}
\newcommand{\AtenXEst}{-12.9}
\newcommand{\NCafes}{189}
\newcommand{\MissRating}{3}
\newcommand{\MissNratings}{2}
\newcommand{\MissAdj}{3}
\newcommand{\MissType}{1}
\newcommand{\MedianN}{446}
\newcommand{\MeanRating}{4.54}
\newcommand{\NTopRated}{69}
\newcommand{\NFew}{22}
\newcommand{\NMedium}{114}
\newcommand{\NMany}{50}
\newcommand{\NIndep}{129}
\newcommand{\NLocal}{48}
\newcommand{\NIntl}{8}
\newcommand{\RVersion}{4.6.1}


% one figure box; shows a visible note if the file is missing
\newcommand{\fig}[2][0.95]{\IfFileExists{../report/figures/#2}{\includegraphics[width=#1\linewidth]{#2}}{\fbox{Missing: \detokenize{#2}}}}

\title[Google ratings of Prague cafés]{Do more Google ratings mean better ratings?}
\subtitle{Evidence from \NCafes{} Prague cafés}
\author{Joseph Muchnick \and Ryan David}
\institute{Data Analytics}
\date{\today}

\begin{document}

\begin{frame}
\titlepage
\end{frame}

\begin{frame}{Question and data}
\begin{itemize}
  \item \textbf{Question:} how is a café's Google rating related to its number of ratings, its type, and a shrinkage-adjusted rating?
  \item \textbf{Data:} student-collected convenience sample, \NCafes{} cafés after cleaning (299 raw rows, 99 empty rows and 9 duplicates removed)
  \item \textbf{Outcome:} Google rating in points (mean \MeanRating{}; median number of ratings \MedianN{})
  \item \textbf{Adjusted rating:} $(nR + 50C)/(n + 50)$ pulls ratings with few votes towards the mean $C$
  \item \textbf{Ten planned questions} (A1--A10), hypotheses written before looking at the data
\end{itemize}
\end{frame}

\begin{frame}{How we answered}
\begin{itemize}
  \item $\alpha = .05$; every result shows an estimate, its 95\% CI and an effect size
  \item Test choice from the descriptives: parametric if the data are normal or the central limit theorem applies (bootstrap check, 5,000 resamples); otherwise rank-based
  \item Regressions: robust (HC3) standard errors when variance is not constant; residual and influence checks
  \item Missing values stay missing: each analysis reports how many cafés it uses
  \item Observational convenience sample: results describe \emph{this sample}, no causal claims
\end{itemize}
\end{frame}

\begin{frame}{A1: Is the mean rating different from 4.5?}
\small\textit{\textbf{Question:} Is the average Google rating of Prague cafés different from 4.5?}\par\medskip
\begin{columns}[T]
\column{0.42\textwidth}\fig{A1_figure.pdf}
\column{0.56\textwidth}
\begin{itemize}
  \item \textbf{Test:} one-sample $t$-test, two-tailed
  \item Mean \AoneEst{} points, \CI{\AoneLo}{\AoneHi} ($n=\AoneN$)
  \item \AoneStat, \AoneP; Cohen's $d=\AoneEff$
  \item \textbf{Takeaway:} Slightly above 4.5, but practically tiny.
\end{itemize}
\end{columns}
\end{frame}

\begin{frame}{A2: More ratings, lower rating? (above vs.\ below median)}
\small\textit{\textbf{Question:} Do cafés with an above-median number of ratings have lower Google ratings than cafés with a below-median number?}\par\medskip
\begin{columns}[T]
\column{0.42\textwidth}\fig{A2_figure.pdf}
\column{0.56\textwidth}
\begin{itemize}
  \item \textbf{Test:} Welch $t$-test, one-tailed (above $<$ below, fixed in advance)
  \item Difference above $-$ below: \AtwoEst{} points, \CI{\AtwoLo}{\AtwoHi} ($n=\AtwoN$)
  \item $d=\AtwoEff$, \CI{\AtwoEffLo}{\AtwoEffHi}; \AtwoStat, \AtwoP
  \item \textbf{Takeaway:} No evidence of a difference.
\end{itemize}
\end{columns}
\end{frame}

\begin{frame}{A3: Adjusted vs.\ raw rating (paired)}
\small\textit{\textbf{Question:} Is the Adjusted Rating different from the raw rating, and if so, by how much?}\par\medskip
\begin{columns}[T]
\column{0.46\textwidth}\fig{A3_figure.pdf}
\column{0.52\textwidth}
\begin{itemize}
  \item \textbf{Test:} paired $t$-test (adjusted $-$ raw)
  \item Mean difference \AthreeEst{} points, \CI{\AthreeLo}{\AthreeHi} ($n=\AthreeN$)
  \item $d_z=\AthreeEff$, \CI{\AthreeEffLo}{\AthreeEffHi}; \AthreeStat, \AthreeP
  \item \textbf{Takeaway:} No average change; a few low-vote cafés move up.
\end{itemize}
\end{columns}
\end{frame}

\begin{frame}{A4: Few, medium or many ratings?}
\small\textit{\textbf{Question:} Does the Google rating differ between cafés with few (under 100), medium (100--1,000) and many (over 1,000) ratings, and if so, how?}\par\medskip
\begin{columns}[T]
\column{0.5\textwidth}\fig{A4_figure.pdf}
\column{0.48\textwidth}
\begin{itemize}
  \item \textbf{Test:} Welch's ANOVA, Games--Howell post hoc
  \item \AfourStat, \AfourP, $\eta^2=\AfourEff$, \CI{\AfourEffLo}{\AfourEffHi} ($n=\AfourN$)
  \item Means: few 4.32, medium 4.62, many 4.48 points
  \item \textbf{Takeaway:} Medium-sized cafés rate highest, not a straight line.
\end{itemize}
\end{columns}
\end{frame}

\begin{frame}{A5: Top-rated (4.7+) by below/above median number of ratings}
\small\textit{\textbf{Question:} Is a café with a below-median number of ratings more likely to be top-rated (4.7 or higher), and if so, by how many percentage points?}\par\medskip
\begin{columns}[T]
\column{0.42\textwidth}\fig{A5_figure.pdf}
\column{0.56\textwidth}
\begin{itemize}
  \item \textbf{Test:} chi-square on the $2\times2$ table, one-tailed
  \item Top-rated: 45.2\% (below median) vs.\ 28.7\% (above)
  \item Difference \AfiveEst{} pp, \CI{\AfiveLo}{\AfiveHi}; Cram\'er's $V=\AfiveEff$, \CI{\AfiveEffLo}{\AfiveEffHi}
  \item \AfiveStat, \AfiveP
  \item \textbf{Takeaway:} Fewer ratings, more often top-rated.
\end{itemize}
\end{columns}
\end{frame}

\begin{frame}{A6: Rating and number of ratings (correlation)}
\small\textit{\textbf{Question:} Is the Google rating correlated with the number of ratings, and if so, in which direction and how strongly?}\par\medskip
\begin{columns}[T]
\column{0.5\textwidth}\fig{A6_figure.pdf}
\column{0.48\textwidth}
\begin{itemize}
  \item \textbf{Test:} Spearman (number of ratings is extremely skewed)
  \item $\rho=\AsixEst$, bootstrap \CI{\AsixLo}{\AsixHi}, $n=\AsixN$; \AsixP
  \item Pearson $r=\AsixPearson$ (\AsixPearsonP); Kendall $\tau=\AsixKendall$ (\AsixKendallP)
  \item \textbf{Takeaway:} At most a weak negative link.
\end{itemize}
\end{columns}
\end{frame}

\begin{frame}{A7: Simple regression: rating per 1,000 ratings}
\small\textit{\textbf{Question:} By how many points does the Google rating change per 1,000 ratings?}\par\medskip
\begin{columns}[T]
\column{0.5\textwidth}\fig{A7_figure.pdf}
\column{0.48\textwidth}
\begin{itemize}
  \item \textbf{Test:} OLS, slope (Breusch--Pagan $p=\AsevenBP$, classical SEs)
  \item Slope \AsevenEst{} points per 1,000 ratings, \CI{\AsevenLo}{\AsevenHi}; \AsevenP; $R^2=\AsevenEff$
  \item \textbf{Takeaway:} No evidence of a slope.
  \item \textbf{Caution:} Cafe Louvre (23,951 ratings) is highly influential. Without it (exploratory): \AsevenXEst, \AsevenXP.
\end{itemize}
\end{columns}
\end{frame}

\begin{frame}{A8: Rating by café type, controlling for number of ratings}
\small\textit{\textbf{Question:} Does the rating differ by Café Type (independent, local chain, international chain) once you account for the number of ratings, and if so, how much of the difference remains?}\par\medskip
\begin{columns}[T]
\column{0.5\textwidth}\fig{A8_figure.pdf}
\column{0.48\textwidth}
\begin{itemize}
  \item \textbf{Test:} multiple regression, joint $F$ test (reference: independent)
  \item \AeightStat, \AeightP; $n=\AeightN$, $R^2=\AeightEff$
  \item International chains: \AeightIntEst{} points, \CI{\AeightIntLo}{\AeightIntHi}
  \item Local chains: \AeightLocEst{} points, \CI{\AeightLocLo}{\AeightLocHi}
  \item \textbf{Takeaway:} International chains rate about 0.29 points lower.
\end{itemize}
\end{columns}
\end{frame}

\begin{frame}{A9: Does the adjustment work differently by number of ratings?}
\small\textit{\textbf{Question:} Does the relationship between raw and Adjusted Rating differ between cafés with few (under 100) and many (100 or more) ratings, and if so, how?}\par\medskip
\begin{columns}[T]
\column{0.5\textwidth}\fig{A9_figure.pdf}
\column{0.48\textwidth}
\begin{itemize}
  \item \textbf{Test:} regression with interaction, robust (HC3) SEs
  \item Slope adjusted on raw: \AnineSlopeU{} (under 100) vs.\ \AnineSlopeO{} (100+)
  \item Interaction \AnineEst, \CI{\AnineLo}{\AnineHi}; \AnineP; $R^2=\AnineEff$
  \item \textbf{Takeaway:} The adjustment pulls few-vote cafés to the mean (by design).
\end{itemize}
\end{columns}
\end{frame}

\begin{frame}{A10: Probability of being top-rated per 1,000 ratings}
\small\textit{\textbf{Question:} Does the probability that a café is top-rated (4.7 or higher) change with each additional 1,000 ratings, and if so, by how many percentage points?}\par\medskip
\begin{columns}[T]
\column{0.5\textwidth}\fig{A10_figure.pdf}
\column{0.48\textwidth}
\begin{itemize}
  \item \textbf{Test:} logistic regression; average marginal effect
  \item \AtenEst{} percentage points per 1,000 ratings, bootstrap \CI{\AtenLo}{\AtenHi}
  \item Odds ratio \AtenOR{} (\CI{\AtenORLo}{\AtenORHi}); \AtenP; $n=\AtenN$ ($\AtenEvents$ top-rated)
  \item \textbf{Takeaway:} More ratings, lower chance of a top rating.
\end{itemize}
\end{columns}
\end{frame}

\begin{frame}{All results at a glance}
\centering\scriptsize
\IfFileExists{../report/tables/summary.tex}{\resizebox{0.95\linewidth}{!}{\begin{tabular}{lrllr}
\toprule
Analysis & n & Estimate & Effect size & $p$ \\
\midrule
A1 one-sample & 186 & 4.54 [4.51, 4.58] & $d$ = 0.17 & .019 \\
A2 two groups & 185 & -0.04 [-0.12, 0.03] & $d$ = -0.17 [-0.46, 0.12] & .123 \\
A3 paired & 186 & 0.010 [-0.003, 0.024] & $d_z$ = 0.11 [-0.03, 0.26] & .125 \\
A4 three groups & 186 & Welch $F$ = 12.05 & $\eta^2$ = 0.15 [0.06, 0.31] & \textless{} .001 \\
A5 two categories & 185 & 15.5 [1.8, 29.3] & $V$ = 0.16 [0.03, 0.30] & .014 \\
A6 correlation & 186 & -0.12 [-0.27, 0.04] & $\rho^2$ = 0.015 & .096 \\
A7 simple regression & 186 & -0.011 [-0.029, 0.006] & $R^2$ = 0.009 & .206 \\
A8 multiple regression & 185 & -0.287 [-0.467, -0.108] & $R^2$ = 0.061 & .007 \\
A9 interaction & 186 & 0.497 [0.401, 0.592] & $R^2$ = 0.983 & \textless{} .001 \\
A10 logistic & 186 & -12.8 [-23.2, -5.6] & OR = 0.56 [0.36, 0.87] & .010 \\
\bottomrule
\end{tabular}
}}{Missing summary table}
\end{frame}

\begin{frame}{Conclusions and limits}
\begin{itemize}
  \item Rating is not a simple function of popularity: cafés with few ratings rate lowest, medium-sized the highest, and there is no clear overall slope (A4, A6, A7)
  \item Cafés with many ratings are less often top-rated (A5, A10); international chains rate lower (A8)
  \item The shrinkage adjustment changes mainly cafés with few votes (A3, A9)
  \item \textbf{Limits:} convenience sample, observational data, cafés in the same district are not independent, A8 rests on \NIntl{} chains, exploratory checks are leads only
  \item Code, report and slides: \texttt{github.com/muchnickjoseph2/cafe-analysis}
\end{itemize}
\end{frame}

\end{document}